% !TeX root = ../main.tex
% Empirical design. Opens Part III: the progression from preliminary design
% by classification, through analytical checking, to detailed numerical
% verification.
\chapter{Empirical Design Methods}
\label{ch:empirical-design}

Design of an underground opening proceeds through methods of increasing
resolution and increasing cost. Empirical methods come first. They convert a
rock mass classification directly into a support recommendation, need only the
data already gathered during characterisation, and can be applied to a whole
alignment in an afternoon. That makes them the natural tool for concept and
preliminary design, for sizing the problem, and for a sanity check on anything
the later methods produce.

They are also the least transferable. An empirical chart encodes the case
history base from which it was built, and says nothing reliable about ground
outside that base --- high stress, squeezing, very weak rock, unusual spans, or
excavation methods that did not exist when the chart was compiled. The
analytical methods of Chapter~\ref{ch:ccm} and the numerical methods of
Chapter~\ref{ch:numerical-modelling} exist to resolve what the charts cannot.

\section{Design stage and method selection}
\label{sec:design-stage}

\subsection{From preliminary to detailed design}
\label{sec:preliminary-to-detailed}

The three method families map onto design stage, and the mapping is driven by
what data exist and what decision is being taken.

Empirical methods suit concept and preliminary design. They need only the
classification data already produced by characterisation, they can be applied
to an entire alignment in a day, and they return the quantity a preliminary
design actually needs --- an indication of support type and quantity, chainage
by chainage. They are also the right tool for screening: identifying which
reaches are routine and which require attention, so that later effort can be
directed at the second group.

Analytical methods suit reference design. They require parameters rather than
ratings, they address one mechanism at a time, and they establish whether the
support proposed will reach equilibrium with the ground and at what
convergence. Their value is diagnostic: they say why a support system works,
which an empirical chart cannot.

Numerical methods suit detailed design and verification. They accommodate real
geometry, staged sequences, interaction between openings, non-uniform ground
and coupled processes. They cost the most, in analyst time and in the data
needed to justify their inputs, and they are the only route to problems the
other two cannot represent at all.

The progression is one of increasing resolution and increasing cost, and it
should be traversed only as far as the decision requires. A support class for
a routine reach in fair ground does not need a finite element model, and
producing one does not make the answer better.

\subsection{What each method can and cannot establish}
\label{sec:method-limits}

Each family returns something different, and confusion about what is being
returned is the source of most misapplication.

Empirical methods return a support class. They do not return a factor of
safety, a convergence, or a load in the support --- and a $Q$-chart support
recommendation carries no statement about how close to failure the excavation
is. What it carries is the accumulated outcome of the case histories behind
it: excavations of this quality and this dimension, supported in this way,
performed acceptably. That is genuine evidence, and it is silent about ground
outside the case base.

Analytical methods return a mechanism and an equilibrium. Given a stress
field, a strength criterion and a support stiffness, the
convergence--confinement method establishes where ground and support reach
equilibrium and how much margin remains. It does not return a construction
sequence, and it cannot represent the geometry of a real excavation: the
circular opening in a hydrostatic field is an idealisation whose departure
from reality must be judged separately.

Numerical methods return a complete field solution --- stresses,
displacements, yield zones everywhere --- and the completeness is what makes
them dangerous. The solution is exact for the model, and its correspondence
with the ground is set entirely by the parameters and constitutive assumptions
supplied. A model reports the same three decimal places whether its modulus
was measured on fifty specimens or assumed from a table.

\subsection{Using the methods together}
\label{sec:methods-together}

Applying all three to the same problem is not duplication, because they fail
in different ways and their agreement is therefore informative.

Where independent methods converge, the result rests on more than one set of
assumptions and can be relied on with correspondingly more confidence. Where
they diverge, the divergence is information rather than an inconvenience: it
usually means that the methods are representing different mechanisms, and
identifying the mechanism each has captured establishes which of them governs.

A common pattern illustrates this. Empirical charts recommend light support;
analytical calculation shows an ample margin; and a numerical model predicts a
large yield zone. The reconciliation is often that the numerical model has
been given a strength-reducing parameter set appropriate to the rock mass but
applied at a scale where structure, not continuum yielding, actually governs
--- or, conversely, that the empirical chart's case base contains nothing at
this stress level. Neither answer is dismissed; the mechanism is identified.
The procedure for doing so systematically is set out in
Chapter~\ref{ch:integrated-design}.

\section{Screening the governing mechanism}
\label{sec:mechanism-screening}

Underground excavations fail in three principal modes, and the governing one
has to be identified before support is selected.
Chapter~\ref{ch:ground-behaviour} classified ground by mechanism and
Chapter~\ref{ch:high-stress} treated the high-stress cases in detail. What
follows reduces that treatment to the screening quantities available at
preliminary stage.

The first screen is the decision tree associated with
\citet{martin1999spalling}. It splits the rock mass on GSI into weak, brittle
and transitional, applies a different strength-to-stress test to each of the
first two, and refers the third to detailed analysis. In weak and soft rock,
with GSI below 30, the test is whether \(\sigma_{ci}/\sigma_v < 2\): if it
is, stress-induced plastic yielding governs, and if not, gravity-induced
movement of structurally defined blocks. In brittle and hard rock, with GSI
above 40, the test is whether \(\sigma_1/\sigma_{ci} > 0.15\): if it is,
stress-induced brittle spalling governs, and if not, structurally controlled
block movement again. Between GSI 30 and 40 lies a transition in which mixed
mechanisms are possible and neither test applies. Here \(\sigma_{ci}\) is the
uniaxial compressive strength of the intact rock, \(\sigma_v\) the vertical
in-situ stress and \(\sigma_1\) the major in-situ stress.
Table~\ref{tab:decision-tree-cases} applies the tree to three cases, and the
three subsections that follow take its branches in turn.

\begin{table}[htbp]
\caption{The decision tree applied to three cases.}
\label{tab:decision-tree-cases}
\centering
\small
\begin{tabular}{lcccc}
\toprule
\textbf{Case} & \textbf{GSI} & $\boldsymbol{\sigma_1/\sigma_{ci}}$ & $\boldsymbol{\sigma_{ci}/\sigma_v}$ & \textbf{Failure type} \\
\midrule
Hard rock mine & 55 & $0.27 > 0.15$ & -- & Stress-induced spalling \\
Weak rock tunnel & 10 & -- & $17.9 > 2$ & Structurally controlled \\
Transition zone & 35 & -- & -- & Neither test applies; detailed analysis \\
\bottomrule
\end{tabular}
\end{table}

\subsection{Stress-induced brittle failure}
\label{sec:screening-brittle}

Stress-induced failure occurs in hard, brittle rock under high stress and
takes the form of spalling, popping or rockburst. It is typical of deep
excavations, at depths of several hundred metres and beyond
(Section~\ref{sec:brittle-failure}).

\paragraph{Depth of spalling.}
The depth of brittle failure follows the empirical relation of
\citet{martin1999spalling},
\[
\frac{r}{a} = 0.5 + 1.25 \frac{\sigma_{\max}}{\mathrm{UCS}},
\]
where $r$ is the radius to the outer limit of the spalled zone and $a$ the
tunnel radius, so that the depth of spalling is $(r/a-1)\,a$ and spalling
initiates at $r/a = 1$, that is at $\sigma_{\max}/\mathrm{UCS} = 0.4$. Here
\[
\sigma_{\max} = 3\sigma_1 - \sigma_3
\]
is the maximum tangential stress at the boundary of a circular opening, from
the Kirsch solution~\citep{kirsch1898}. In this expression \(\sigma_1\) and
\(\sigma_3\) are the principal stress components in the plane normal to the
tunnel axis. The principal stresses of the in-situ field may be entered
directly only where the axis is parallel to one of them.

The relation applies only to brittle spalling in massive hard rock. It
assumes a circular opening in an elastic stress field, and it is not to be
used for weak, highly fractured or anisotropic rock masses.

\paragraph{Severity.}
The crack initiation stress CI, typically 0.4 to 0.6 of the uniaxial
compressive strength, gives a second measure. The ratio
\(\sigma_{\max}/\mathrm{CI}\) indicates the severity of failure: below 1 the
opening is stable, with no stress damage; between 1 and 1.5 minor spalling is
expected; between 1.5 and 2 moderate overbreak; and above 2 serious
overbreak. Table~\ref{tab:spalling-example} works both measures through for a
strong rock under a major stress of 60~MPa. The two bounds on CI place the
same opening in adjacent severity classes, which is a fair statement of the
resolution of the method.

\begin{table}[htbp]
\caption{Illustrative calculation of spalling depth and severity.}
\label{tab:spalling-example}
\centering
\small
\begin{tabular}{lcl}
\toprule
\textbf{Parameter} & \textbf{Value} & \textbf{Notes} \\
\midrule
UCS & 224~MPa & Laboratory testing \\
$\sigma_1$ & 60~MPa & In-situ stress \\
$\sigma_3$ & 11~MPa & \\
$\sigma_{\max}$ & $3\times 60 - 11 = 169$~MPa & Maximum tangential stress \\
CI (low) & $0.4\times224 = 89.6$~MPa & \\
CI (high) & $0.6\times224 = 134.4$~MPa & \\
$\sigma_{\max}/\mathrm{CI}$ & 1.89 (low), 1.26 (high) & Moderate to minor spalling \\
$r/a$ & $0.5 + 1.25\times169/224 = 1.44$ & Spalling depth = 44\% of radius \\
\bottomrule
\end{tabular}
\end{table}

\subsection{Structurally controlled failure}
\label{sec:screening-structural}

Structurally controlled failure is the gravity-driven fall or sliding of
blocks defined by intersecting joint sets. It occurs where the rock mass is
blocky and the stress is low relative to strength, and its analysis is the
kinematic one of Chapter~\ref{ch:structural-analysis}.

This is the mode for which the empirical systems were built. RMR and Q are
most reliable in preliminary support design where behaviour is governed
mainly by rock mass quality and structure. Severe brittle failure, high
stress and squeezing lie outside their case base and require the analysis of
Chapter~\ref{ch:high-stress}; dynamic rockburst conditions in particular
cannot be addressed by a classification rating alone.

\subsection{Squeezing}
\label{sec:screening-squeezing}

Squeezing, the development of large deformation, occurs in weak rock masses
of poor quality under high in-situ stress. It is common in fault zones and
sheared rock and in tunnelling at great depth through weak strata
(Section~\ref{sec:squeezing-swelling}).

The strain relation of \citet{hoekmarinos2000squeezing} gives a first
estimate of its severity,
\[
\varepsilon = 0.2\left(\frac{\sigma_{cm}}{p_0}\right)^{-2} ,
\]
where $\varepsilon$ is the closure strain of the unsupported tunnel in per
cent, \(\sigma_{cm}\) is the strength of the rock mass and
\[
p_0 = \frac{\sigma_1+\sigma_2+\sigma_3}{3}
\]
is the mean in-situ stress. The rock mass strength in this relation is the
one its authors defined,
\[
\sigma_{cm} = \left(0.0034\,m_i^{0.8}\right)\sigma_{ci}
\left[1.029 + 0.025\,e^{-0.1 m_i}\right]^{\mathrm{GSI}} ,
\]
a strength averaged over the confinement around a tunnel. It exceeds the
uniaxial value \(\sigma_{ci}\,s^{a_{HB}}\) of
Chapter~\ref{ch:rock-mass-parameters}, by a factor of about three in poor
rock, and because the strain varies as the inverse square of the strength,
substituting the uniaxial value there overstates the strain roughly
tenfold. The strain expression is an approximate curve fit, suitable for
preliminary assessment only.
Table~\ref{tab:squeezing-classification} gives the classes attached to the
strain. With $\sigma_{cm} = 10$~MPa and $p_0 = 25$~MPa, for example,
\[
\frac{\sigma_{cm}}{p_0}=0.4 ,
\qquad
\varepsilon = 0.2(0.4)^{-2} = 1.25\% ,
\]
which corresponds to minor squeezing.

\begin{table}[htbp]
\caption{Squeezing classification, after \citet{hoekmarinos2000squeezing}.}
\label{tab:squeezing-classification}
\centering
\small
\begin{tabular}{ll}
\toprule
\textbf{Strain $\varepsilon$} & \textbf{Classification} \\
\midrule
$<1\%$ & Few support problems \\
$1$--$2.5\%$ & Minor squeezing \\
$2.5$--$5\%$ & Severe squeezing \\
$5$--$10\%$ & Very severe squeezing \\
$>10\%$ & Extreme squeezing \\
\bottomrule
\end{tabular}
\end{table}

\section{Empirical support design}
\label{sec:empirical-support-design}

Both classification systems lead to a support recommendation through a chart
compiled from case histories. The charts themselves are in the sources
cited; the tables of this section carry the quantities read from them.

\subsection{The RMR system}
\label{sec:rmr-support-design}

\paragraph{Unsupported span and stand-up time.}
The chart of \citet{bieniawski1989rmr} relates RMR to the maximum
unsupported roof span and to the stand-up time. Two bounding curves divide it
into three zones. Below the lower bound the span requires no support, and
above the upper bound it collapses immediately. Between the bounds the roof
fails in time: it stands for a finite period that the chart returns, and the
design question is whether the support cycle can be completed within that
period. This intermediate zone is the condition of most excavations. The
chart is entered at the RMR value, and the intersections with the two bounds
give the maximum unsupported span and the corresponding stand-up time.
Table~\ref{tab:rmr-span-standup} summarises the readings by class.

\begin{table}[htbp]
\caption{Approximate unsupported spans and stand-up times read from the chart of \citet{bieniawski1989rmr}. The values are indicative; a design reading is taken from the chart itself.}
\label{tab:rmr-span-standup}
\centering
\small
\begin{tabularx}{\textwidth}{l c X}
\toprule
\textbf{RMR} & \textbf{Approx.\ max span (m)} & \textbf{Approx.\ stand-up time} \\
\midrule
81--100 (Class I) & 15--30 & $>10$ years \\
61--80 (Class II) & 10--15 & Months to years \\
41--60 (Class III) & 5--10 & Weeks to months \\
21--40 (Class IV) & 2.5--5 & Hours to days \\
$<20$ (Class V) & $<2.5$ & Minutes (immediate collapse) \\
\bottomrule
\end{tabularx}
\end{table}

\paragraph{Adjustment for joint orientation.}
The basic RMR is corrected for the orientation of the discontinuities
relative to the structure by the amounts in
Table~\ref{tab:rmr-orientation-adjustment},
\[
\mathrm{Adjusted\ RMR} = \mathrm{RMR} + \mathrm{Adjustment} ,
\]
and the most unfavourable joint set governs.

\begin{table}[htbp]
\caption{RMR adjustment for joint orientation~\citep{bieniawski1989rmr}. The most unfavourable joint set governs.}
\label{tab:rmr-orientation-adjustment}
\centering
\small
\begin{tabular}{lccccc}
\toprule
\textbf{Condition} & \textbf{Very favourable} & \textbf{Favourable} & \textbf{Fair} & \textbf{Unfavourable} & \textbf{Very unfavourable} \\
\midrule
Tunnels & 0 & $-2$ & $-5$ & $-10$ & $-12$ \\
Foundations & 0 & $-2$ & $-7$ & $-15$ & $-25$ \\
\bottomrule
\end{tabular}
\end{table}

\paragraph{Support recommendations.}
Table~\ref{tab:rmr-support} gives the excavation method and support
recommended for each class. The recommendations are written for a span of
10~m. For a markedly different span the table indicates a class of support
and not its quantities, and of those quantities only the bolt length is
scaled. Two relations are in use for a span $B$ between 6 and 18~m,
\[
L = 0.125B + 2.25
\qquad \text{or} \qquad
L = 0.3B ,
\]
with \(L\) and \(B\) in metres. They agree near a span of 13~m and diverge
towards the ends of the range, the first giving the longer bolt by two-thirds
at 6~m and the second by a fifth at 18~m. Both are therefore computed and
the longer adopted, unless the shorter can be justified against local
experience.

\begin{table}[htbp]
\caption{RMR support recommendations for a 10~m span~\citep{bieniawski1989rmr}.}
\label{tab:rmr-support}
\centering
\small
\begin{tabularx}{\textwidth}{l l X}
\toprule
\textbf{RMR} & \textbf{Excavation method} & \textbf{Support} \\
\midrule
81--100 & Full face, 3~m advance & Generally no support, except spot bolting where needed \\
61--80 & Full face, 1--1.5~m advance & Bolts 3~m @ 2.5~m in crown; 50~mm shotcrete in crown \\
41--60 & Top heading and bench & Systematic bolts 4~m @ 1.5--2~m; 50--100~mm shotcrete \\
21--40 & Top heading and bench & Systematic bolts 4--5~m @ 1--1.5~m; 100--150~mm shotcrete; light ribs \\
$<20$ & Multiple drifts & Systematic bolts 5--6~m @ 1--1.5~m; 150--200~mm shotcrete; heavy ribs \\
\bottomrule
\end{tabularx}
\end{table}

\subsection{The Q-system}
\label{sec:q-support-design}

\paragraph{Equivalent dimension.}
The Q-system support chart~\citep{barton1974q,barton1994qsystem} is entered
not with the physical span but with an equivalent dimension,
\[
D_e = \frac{\text{span or height}}{\text{ESR}} ,
\]
the span or wall height in metres divided by the excavation support ratio
(ESR). The ratio encodes the consequence of failure
(Table~\ref{tab:esr}). It is unity for a major tunnel and below unity where
the opening is permanent and heavily used, so that a more important
excavation presents a larger equivalent dimension and attracts heavier
support.

\begin{table}[htbp]
\caption{Values of the excavation support ratio (ESR).}
\label{tab:esr}
\centering
\small
\begin{tabularx}{\textwidth}{X c}
\toprule
\textbf{Excavation type} & \textbf{ESR} \\
\midrule
Temporary mine openings & 3--5 \\
Permanent mine openings, water tunnels & 1.6 \\
Minor road/rail tunnels, surge chambers & 1.3 \\
Power stations, major tunnels, civil defence & 1.0 \\
Underground railway stations, sports facilities & 0.8 \\
Very important caverns, long lifetime ($>100$~y) & 0.5 \\
\bottomrule
\end{tabularx}
\end{table}

\paragraph{Reading the chart.}
The chart plots $Q$ on the horizontal axis and $D_e$ on the left vertical
axis. The intersection of the two fixes a support category, numbered 1 to 9
(Table~\ref{tab:q-support-categories}), and with it the recommended thickness
of fibre-reinforced sprayed concrete. From the same point the average bolt
spacing is read upward for areas with fibre-reinforced sprayed concrete and
downward for areas without it. The bolt length is read on the right vertical
axis, which is drawn for an ESR of 1; for any other ESR the length is
therefore read against the physical span or height and not against the
equivalent dimension. In the heavier categories the sprayed concrete is
supplemented by reinforced ribs of sprayed concrete (RRS), graded I and II in
order of increasing section.

\begin{table}[htbp]
\caption{Simplified Q-system support categories. RRS denotes reinforced ribs of sprayed concrete.}
\label{tab:q-support-categories}
\centering
\small
\begin{tabular}{ll}
\toprule
\textbf{Category} & \textbf{Support description} \\
\midrule
1 & Unsupported or spot bolting \\
2 & Spot bolting \\
3 & Systematic bolting, fibre-reinforced shotcrete 5--6~cm \\
4 & Fibre-reinforced shotcrete 6--9~cm + bolting \\
5 & Fibre-reinforced shotcrete 9--12~cm + bolting \\
6 & Shotcrete 12--15~cm + RRS~I + bolting \\
7 & Shotcrete $>15$~cm + RRS~II + bolting \\
8 & Cast concrete lining \\
9 & Special evaluation \\
\bottomrule
\end{tabular}
\end{table}

\paragraph{Walls and the unsupported span.}
Walls are designed on the same chart with an adjusted rating: $Q_{\text{wall}}
= 5Q$ where $Q > 10$, $Q_{\text{wall}} = 2.5Q$ where $0.1 < Q < 10$, and the
actual $Q$ where $Q < 0.1$. The chart is then entered with
\[
D_e = \frac{\text{wall height}}{\text{ESR}} .
\]
The maximum unsupported span follows from the lower limit of the chart,
\[
\mathrm{Span}_{\max} = 2 \, \mathrm{ESR} \, Q^{0.4} ,
\]
in metres.

\subsection{A station cavern by both systems}
\label{sec:station-cavern-example}

A railway station cavern has a span of 16~m and a wall height of 6~m in rock
of RMR = 50 and $Q = 5$. Two joint sets are present. The first strikes
perpendicular to the cavern axis and dips at $30^\circ$; the second strikes
parallel to the axis and dips at $70^\circ$.

\paragraph{RMR design.}
The first joint set, with the drive taken in the direction of dip, is
favourable ($-2$), and the second is very unfavourable ($-12$). The worst
case governs, so the adjusted RMR is $50 - 12 = 38$, which is Class~IV. The
bolt length from the two relations is
\[
L = 0.125\times16 + 2.25 = 4.25~\text{m}
\qquad \text{or} \qquad
L = 0.3\times16 = 4.8~\text{m} ,
\]
and the longer is adopted. The design is therefore 4.8~m bolts at 1~m
spacing with 100--150~mm of shotcrete and light ribs. Because the support
table is written for a 10~m span, it supplies the class of support for this
16~m opening and not its quantities, the bolt length excepted.

\paragraph{Q-system design.}
A station takes ESR = 0.8, so for the roof $D_e = 16/0.8 = 20$. The chart
gives support category 5, fibre-reinforced shotcrete of 9--12~cm with
bolting, at a bolt spacing of roughly 2.1~m. The bolt length, read against
the physical span of 16~m because the ESR is not unity, is about 4~m. For the
walls, $Q_{\text{wall}} = 2.5\times 5 = 12.5$ and $D_e = 6/0.8 = 7.5$, which
gives category 3, systematic bolting with 5--6~cm of shotcrete.

The maximum unsupported span for $Q = 5$ and ESR = 0.8 is
$2(0.8)(5^{0.4}) = 3.0$~m, against a designed span of 16~m, and the RMR route
gives 2.5--5~m for the adjusted Class~IV. The two systems agree,
independently, that nothing in this cavern stands unsupported.

\paragraph{The divergence between the two.}
The RMR design (4.8~m bolts, 100--150~mm of shotcrete and light ribs) is more
conservative than the Q-system design (4~m bolts and 9--12~cm of shotcrete).
Part of the divergence is structural. The RMR route carries an explicit
penalty of 12 points for the steeply dipping set striking parallel to the
axis, whereas $Q$ has no orientation term, so the same adverse geometry is
invisible to it. The Q recommendation is the lighter of the two by
construction, and the effect of orientation has to be imposed on it
separately. The remaining choice is a matter of engineering judgement, which
weighs the uncertainty of the ratings, the possibility of mixed failure
modes, the importance of the cavern (ESR = 0.8) and local experience.
Empirical systems are guides and not prescriptions.

\section{Calibration and the influence of uncertainty}
\label{sec:case-study-and-the-influence-of-uncertainty}

\paragraph{A shallow car park cavern in sandstone.}
The underground car park at the Sydney Opera House shows both systems at
their limits~\citep{pells2002developments}. RMR discriminated poorly,
identifying only three classes. $Q$ ranged from 43 to 0.009 and so captured
the variability better, yet it predicted lighter support than was actually
installed. Empirical systems have to be calibrated to local conditions.
Neither RMR nor Q is universally applicable, and the reliability of each
depends on how closely the case histories behind it match the project in
hand.

\paragraph{Uncertainty of the inputs.}
The empirical systems rest on parameters that carry inherent uncertainty
(Chapter~\ref{ch:site-investigation}). RQD is subject to orientation bias
and core loss. Joint spacing and joint condition depend on the
interpretation made during logging. Groundwater is often inferred from
limited measurements. The stress reduction factor (SRF) relies on an estimate
of the in-situ stress. Because the inputs are uncertain, the support
recommendation is uncertain also.

\paragraph{Effect on the recommendation.}
The effect depends on where the rating lies. If RMR = $55 \pm 5$, both
bounds remain in Class~III (41--60) and the recommended support class does
not change; the uncertainty surfaces instead in the quantity selected within
the class. A rating near a class boundary behaves differently: at RMR =
$42 \pm 5$ the lower bound crosses into Class~IV and the recommendation
changes discontinuously. Likewise, if $Q = 5 \pm 2$, the lower bound
($Q = 3$) may shift into a heavier support category, and small variations in
joint logging, such as the interpretation of $J_r$, can affect the category.

A useful check is to recompute the support recommendation with the RMR and
$Q$ inputs perturbed by $\pm 10$\,per cent. The comparison identifies which
support elements are governed by a rating that is itself poorly constrained,
and therefore which additional observation would most reduce the uncertainty
carried into the design.
